% all_N_assembly.tex  --  Step-1's contribution to the FULL (all-N) #23 paper (v3).
% Gated on Step-2's exact band bound delta=0. Self-contained section; \input into the
% full paper once delta=0 is in hand. H2 peeling is RETIRED by Lemma~\ref{lem:transfer}
% (GPT-Pro audited 2026-06-26; no-fractional-gap numerically verified by
% verify_graphon_transfer.py: 11 triangle-free graphs incl. C5[2],C5[3],Petersen,Grotzsch).
\section{From the graphon bound to all finite \texorpdfstring{$N$}{N}}\label{sec:allN}

Throughout, $G$ is a finite triangle-free graph on $N$ vertices, $e(G)$ its edge
count, $\mathrm{mc}(G)$ its MaxCut, and $\beta(G)=e(G)-\mathrm{mc}(G)$ its
bipartization number. Let $W_G$ be the $\{0,1\}$ \emph{step graphon} of $G$: split
$[0,1]$ into cells $I_1,\dots,I_N$ of measure $1/N$ and set $W_G\equiv 1$ on
$I_i\times I_j$ and $I_j\times I_i$ whenever $ij\in E(G)$, $W_G\equiv 0$ otherwise.
For a graphon $W$ write
\[
  d_{\mathrm{edge}}(W)=\int_0^1\!\!\int_0^1 W,\qquad
  d_{\mathrm{mono}}(W)=\inf_{S\subseteq[0,1]\ \text{meas.}}
      \Big(\int_{S\times S}W+\int_{S^c\times S^c}W\Big).
\]

\begin{lemma}[Exact graphon transfer]\label{lem:transfer}
For every finite graph $G$ on $N$ vertices,
\[
  d_{\mathrm{edge}}(W_G)=\frac{2e(G)}{N^2},\qquad
  d_{\mathrm{mono}}(W_G)=\frac{2\beta(G)}{N^2}.
\]
Both are exact equalities: there is no $o(1)$, no $O(1/N)$, and no fractional--MaxCut
gap.
\end{lemma}

\begin{proof}
The edge identity is immediate ($e(G)$ blocks of area $2/N^2$). For the second, fix a
measurable $S$ and put $x_i=N\,\mu(S\cap I_i)\in[0,1]$, the fraction of cell $i$ inside
$S$. Since $W_G$ is constant on each block $I_i\times I_j$,
\[
  \int_{S\times S}W_G+\int_{S^c\times S^c}W_G
  =\frac{2}{N^2}\sum_{ij\in E(G)}\big[x_ix_j+(1-x_i)(1-x_j)\big]
  =:\frac{2}{N^2}\,f(x).
\]
The function $f$ is \emph{multilinear} (affine in each $x_i$ separately: the
$x_i$--terms are $\sum_{j\sim i}\big[2x_ix_j-x_i-x_j+1\big]$, whose $x_i$--coefficient
$\sum_{j\sim i}(2x_j-1)$ does not involve $x_i$). A multilinear function on the box
$[0,1]^N$ attains its minimum at a vertex $x\in\{0,1\}^N$. Such a vertex is the
indicator of an integer $2$--colouring, and $f(x)$ counts its monochromatic edges, so
$\min_{x\in[0,1]^N}f(x)=\min_{x\in\{0,1\}^N}f(x)=e(G)-\mathrm{mc}(G)=\beta(G)$.
Therefore $d_{\mathrm{mono}}(W_G)=2\beta(G)/N^2$.
\end{proof}

\begin{theorem}[All-$N$ reduction]\label{thm:allN}
Suppose the exact graphon bound
\[
  d_{\mathrm{mono}}(W)\le\frac{2}{25}\qquad\text{for every triangle-free graphon }W
  \tag{$\delta{=}0$}
\]
holds. Then $\beta(G)\le N^2/25$ for every triangle-free graph $G$ on every $N$, i.e.
Erd\H{o}s problem \#23 holds in full.
\end{theorem}

\begin{proof}
$G$ triangle-free $\Rightarrow$ $W_G$ triangle-free, so $(\delta{=}0)$ applies to
$W_G$: $d_{\mathrm{mono}}(W_G)\le 2/25$. By Lemma~\ref{lem:transfer},
$2\beta(G)/N^2\le 2/25$, hence $\beta(G)\le N^2/25$.
\end{proof}

\noindent\textbf{Assembling $(\delta{=}0)$.} The hypothesis is split by edge density
$d:=d_{\mathrm{edge}}(W_G)=2e(G)/N^2$ into three parts that tile $[0,\tfrac12]$ with
\emph{shared endpoints}:
\[
  d<0.2486,\qquad 0.2486\le d\le 0.3197,\qquad d>0.3197 .
\]
The \emph{closed band} owns both endpoints; this is deliberate, so that no graph falls
in a gap and the low-tail endpoint issue caused by the Balogh--Clemen--Lid\'ick\'y
(BCL) finite-density normalisation $e/\binom{N}{2}$ (vs.\ our $2e/N^2$) is avoided.
\begin{itemize}
\item \emph{Band} $0.2486\le d\le 0.3197$: the exact band bound
  $d_{\mathrm{mono}}(W)\le 2/25$, the sole hard input, supplied by Step-2.
\item \emph{Tails} $d<0.2486$ or $d>0.3197$ \emph{(strict)}: the Balogh--Clemen--Lid\'ick\'y
  theorem (\texttt{arXiv:2103.14179}, Thm.\ on triangle-free bipartisation), whose parts
  \textbf{(b)} $\beta(G)\le n^2/25$ when $e(G)\ge 0.3197\binom n2$ and \textbf{(c)}
  $\beta(G)\le n^2/25$ when $e(G)\le 0.2486\binom n2$ hold \emph{for $n$ large enough}
  (here $\beta=D_2$, the bipartisation number). \emph{The two BCL thresholds
  $0.2486,0.3197$ are exactly our band edges}; we discharge the tails at fixed $N$ by the
  blow-up $\beta(G[t])=t^2\beta(G)$. BCL's hypothesis is on the \emph{finite} density
  $e(G[t])/\binom{Nt}{2}=d\cdot\frac{Nt}{Nt-1}\to d$ (Remark below) -- \emph{not} the
  graphon density -- which for large $t$ lies strictly inside the same open tail as $d$.
  Since BCL parts (b),(c) carry the \emph{exact} constant $n^2/25$ for $n\ge n_0$, a
  \emph{single} $t$ with $Nt\ge n_0$ already gives $\beta(G[t])\le (Nt)^2/25$; dividing
  by $t^2$ yields $\beta(G)\le N^2/25$ \emph{with no limit and no $o(1)$}. (Should one
  prefer to read BCL only asymptotically as $(\tfrac1{25}+o(1))n^2$, the same conclusion
  follows by letting $t\to\infty$ with $N$ fixed; the argument is robust to either reading.)
\end{itemize}

\noindent\emph{(What is assembled.)} The three cases together establish
$d_{\mathrm{mono}}(W_G)\le 2/25$ for \emph{every} finite triangle-free $G$ -- precisely
the instances of $(\delta{=}0)$ that Theorem~\ref{thm:allN} invokes (it applies the
hypothesis only to step graphons $W_G$). The band input holds for arbitrary triangle-free
graphons; the tails are discharged through the finite blow-up $\beta(G[t])=t^2\beta(G)$ and
so for $W_G$. This $W_G$ form is all Theorem~\ref{thm:allN} uses, so no statement about
non-step triangle-free graphons in the tails is needed.

\noindent\textbf{Remark (normalisation at the band edges).} BCL's density is the
graph density $e/\binom{N}{2}=2e/(N(N-1))$, whereas our partition uses the graphon
density $d=2e(G)/N^2$; the two differ by the factor $N/(N-1)\to1$. The blow-up transfer
reconciles them \emph{exactly}: for $G[t]$ on $Nt$ vertices,
\[
  \frac{e(G[t])}{\binom{Nt}{2}}
  =\frac{2\,t^2 e(G)}{Nt(Nt-1)}
  =\frac{2e(G)}{N^2}\cdot\frac{Nt}{Nt-1}
  \ \xrightarrow[t\to\infty]{}\ \frac{2e(G)}{N^2}=d,
\]
and for every finite $t$ this BCL density strictly exceeds $d$ (as $Nt-1<Nt$). Hence:
if $d>0.3197$ then $e(G[t])/\binom{Nt}{2}>d>0.3197$ for \emph{all} $t$; if $d<0.2486$
then $e(G[t])/\binom{Nt}{2}\downarrow d<0.2486$, so it is $<0.2486$ for all large $t$.
Either way arbitrarily large $t$ (so $n=Nt$ is ``large enough'' for BCL) places $G[t]$
strictly inside the asserted open BCL tail, giving
$\beta(G[t])\le\bigl(\tfrac1{25}+o_t(1)\bigr)(Nt)^2$; dividing by $t^2$ and letting
$t\to\infty$ yields $\beta(G)\le N^2/25$. (This makes the tail argument robust to whether
BCL's bound carries the exact constant at finite $n$ or only asymptotically.) Because the \emph{closed} band $[0.2486,0.3197]$ owns both
endpoints (handled by Step-2's bound, not BCL), no graph with $d$ exactly $0.2486$ or
$0.3197$ relies on the tails, and no graph falls in a gap.

\noindent\textbf{Remark (no rounding for $5\nmid N$).} When $5\nmid N$, $N^2/25$ is not
an integer while $\beta(G)$ is; but the conclusion is the \emph{real} inequality
$\beta(G)\le N^2/25$, which at once gives $\beta(G)\le\lfloor N^2/25\rfloor$. No
integrality argument is needed. The integrality rounding of the finite paper (the
threshold $\tfrac{25}{2}N^2\delta<1$) was required \emph{only} for the approximate band
bound $\delta>0$; at $\delta=0$ it is vacuous. In particular the H2 peeling lemma is
\emph{not} used anywhere in Theorem~\ref{thm:allN}.

\noindent\textbf{Remark (a genuine minimum; no ``$N$ large'').} The infimum defining
$d_{\mathrm{mono}}(W_G)$ is attained: $f$ is continuous on the compact box $[0,1]^N$, so
``$\inf$'' is a minimum (realised at the integer cut). No step invokes ``$N$ large
enough'': Lemma~\ref{lem:transfer} is exact at every finite $N$, and the lone asymptotic
input (BCL in the tails) is discharged by the blow-up $t\to\infty$ with $N$ \emph{fixed},
not by enlarging $N$.

\noindent\textbf{Verification.} The no-fractional-gap core of Lemma~\ref{lem:transfer}
(that the continuous minimum of $f$ over $[0,1]^N$ equals the integer $\beta(G)$) is
checked exhaustively-by-restart and by random interior sampling in
\texttt{verify\_graphon\_transfer.py} on C5, $C_5[2]$, $C_5[3]$, the Petersen graph, the
Gr\"otzsch graph, $C_7$, and random triangle-free graphs: in every case the vertex
minimum equals $\beta(G)$ and no interior point undercuts it.
